\thispagestyle{empty}
\begin{center}
\begin{LARGE}
\textbf{Abstract}
\end{LARGE}
\end{center}
\begin{quotation}

This research project aims to develop an automated attendance system based on facial recognition for the Polytechnic Faculty of the National University of the East (FPUNE). Currently, attendance is recorded manually, which often results in human errors and delays in administrative processes.

The proposed solution applies computer vision techniques to detect and recognize faces in real time, automating the attendance process. Tools such as Python, OpenCV, TensorFlow, and models like YOLOv8 and ArcFace were used. The methodology included data collection, model training, and validation in a real environment.

The results demonstrate high accuracy and efficiency, even under variable lighting and viewing conditions. Overall, the system improves academic management by reducing manual workload, minimizing errors, and optimizing teachers’ time during attendance control.

\vspace*{0.5cm}
\noindent \textbf{Key words:} 1. Facial Recognition, 2. Attendance Control, 3. Computer Vision.

\end{quotation}
